\section{Discussion and open questions}
\label{sec:discussion}

The expected-length bounds in \cref{obj:thm:sharp-length-frontier-full-elbow} provide precision benchmarks for the declared source--target experiment: independent samples with \(n\) source records and \(n\) target covariate records, under the restrictions in \cref{obj:def:model-class}. For \(n\ge n_0\), fixed calibration and regularity parameters, and an evaluation floor \(0<a\le1/4\), the two regimes describe the interaction between transported compliance and the rough estimation scale. When \(a>n^{-1/3}\), the optimal worst-case expected length has order \(n^{-1/3}/a\). When \(a\le n^{-1/3}\), it has constant order. Thus increasing the common sample size improves the precision benchmark once the strength floor exceeds the estimation scale, while global honesty requires positive worst-case expected length across the weaker-strength slices.

The floor \(a\) indexes the performance criterion in \cref{obj:def:strength-slice,obj:def:length-frontier}; the interval in \cref{obj:def:score-interval} uses the observed transported contrasts. Its law-specific expected-length guarantee in \cref{obj:thm:honest-upper-full-elbow} relates precision to the actual transported first stage \(\mu(P)\). A single interval therefore supplies global coverage and responds to the compliance strength of the underlying law, with inverse-strength control of expected length above the rough estimation scale and bounded length throughout the effect domain. This guarantee concerns expected length under repeated sampling; it is compatible with variation in the reported interval across realizations.

The fixed value \(\beta=1/8\) specifies the rough-nuisance regularity class analyzed here. At this regularity, the cubic correction is used to control the nonlinear remainder at the required scale. The split-sample estimator includes linear, quadratic, and cubic Taylor corrections across all seven learned marked-density coordinates and has mean-square error bounded by \(C_{\mathrm{mse}}n^{-2/3}\) under the model assumptions for \(n\ge256\). The resulting contrast-estimation radius is therefore of order \(n^{-1/3}\), which becomes the elbow in the length frontier.

The displayed finite-sample radius, \(2\sqrt{2C_{\mathrm{mse}}/\alpha}\,n^{-1/3}\), makes the hidden rate constant explicit: it depends on the density envelopes, the common H\"older radius, and the desired coverage level through \(C_{\mathrm{mse}}\) and \(\alpha\). The threshold \(n_0=256\) ensures that the deterministic four-block construction and its resolution choices are defined. Numerical precision is governed by the calibration constant and the transported first stage. For example, take the admissible inputs \(c_f=0.9\), \(C_f=1.1\), \(L=1.01\), and \(\alpha=0.05\). Substitution into the displayed formula in \cref{obj:lem:marked-cubic-mse} gives \(C_{\mathrm{mse}}\approx1.16\times10^{27}\), so the certified score tolerance is approximately \(4.31\times10^{14}n^{-1/3}\): about \(6.79\times10^{13}\) at \(n=256\) and \(4.31\times10^{12}\) at \(n=10^6\). These values exceed the effect-domain diameter and show that the explicit worst-case envelope is highly conservative; a realized affine-score interval can be shorter, but that realized length is not the uniform certificate. An analyst supplies fixed, scientifically defensible bounds on density overlap and Hölder size before applying this finite-sample calibration. The expected-length bound then uses the corresponding \(C_0\) and the law-specific first stage.

An application of this benchmark requires the source encouragement experiment to satisfy conditional randomization, exclusion, and monotonicity, together with the sampling and regularity conditions of the model. Transport to the target population rests on the conditional equalities for the complier-weighted outcome contrast and the complier share in \cref{obj:ass:transport-complier-outcome,obj:ass:transport-complier-share}. Under these conditions, the precision benchmark concerns the target-population average treatment effect among compliers, using source receipt and outcome observations together with target covariate observations.

Separate recruitment costs motivate the unequal-sample extension of the same decision problem in \cref{obj:def:unequal-sample-handle}. Let \(N_{\mathrm S}\) denote the source sample size and \(N_{\mathrm T}\) the target sample size. Retaining the law-level restrictions while allowing these counts to differ gives the following criterion.

\begin{definitionv}{def:unequal-sample-handle}[Unequal-sample frontier \(R_{N_{\mathrm S},N_{\mathrm T}}(a)\)]
The unequal-sample minimax expected-length frontier is
\[
R_{N_{\mathrm S},N_{\mathrm T}}(a)
=
\inf_{C_{N_{\mathrm S},N_{\mathrm T}}\in
\mathfrak H^{\mathrm u}_{N_{\mathrm S},N_{\mathrm T}}}
\sup_{P\in\mathcal M^{\mathrm u}_{N_{\mathrm S},N_{\mathrm T}}(a)}
\mathbb E_P\!\left[
\lambda_{\Theta}\!\left(C_{N_{\mathrm S},N_{\mathrm T}}\right)
\right],
\qquad
0<a\le\frac14,
\]
for positive integers \(N_{\mathrm S}\) and \(N_{\mathrm T}\), where \(\lambda_{\Theta}\) denotes Lebesgue length restricted to the effect domain \(\Theta\).

Here \(\mathcal M^{\mathrm u}_{N_{\mathrm S},N_{\mathrm T}}\) is the class of full-data and one-record observed-law triples satisfying every law-level density, smoothness, overlap, causal, transport, and positive-strength condition of the equal-sample model class \(\mathcal M_n\), observed under
\[
P_{\mathrm S}^{\otimes N_{\mathrm S}}
\otimes
P_{\mathrm T}^{\otimes N_{\mathrm T}},
\]
where \(P_{\mathrm S}\) is the law of one observed source record and \(P_{\mathrm T}\) is the law of one observed target covariate. The independent source and target samples contain \(N_{\mathrm S}\) records \((X,Z,D,Y)\) and \(N_{\mathrm T}\) records \(X\), respectively. With \(\mu(P)\) denoting the transported first stage, the strength slice is
\[
\mathcal M^{\mathrm u}_{N_{\mathrm S},N_{\mathrm T}}(a)
=
\left\{
P\in\mathcal M^{\mathrm u}_{N_{\mathrm S},N_{\mathrm T}}:
a\le\mu(P)\le\frac14
\right\}.
\]

Let \(C_{N_{\mathrm S},N_{\mathrm T}}\) range over measurable connected random intervals contained in \(\Theta\) and based on this product experiment. With \(\theta(P)\) denoting the target-population average treatment effect among compliers and \(\alpha\) the prescribed noncoverage probability, define
\[
\mathfrak H^{\mathrm u}_{N_{\mathrm S},N_{\mathrm T}}
=
\left\{
C_{N_{\mathrm S},N_{\mathrm T}}:
\inf_{P\in\mathcal M^{\mathrm u}_{N_{\mathrm S},N_{\mathrm T}}}
P\!\left(
\theta(P)\in C_{N_{\mathrm S},N_{\mathrm T}}
\right)
\ge1-\alpha
\right\}.
\]
\end{definitionv}

The unequal-sample minimax expected-length criterion \(R_{N_{\mathrm S},N_{\mathrm T}}(a)\) preserves the distinction between global coverage and performance on a strength slice. Coverage is required over the entire unequal-sample model, while the supremum defining expected length evaluates laws whose transported first stage lies between the specified floor and \(1/4\), inclusive. The two sample counts consequently enter through the observation experiment, with the population restrictions and the estimand held fixed.

\subsection{Open questions}

The unequal-sample criterion makes allocation between source outcomes and target covariates a statistical design question. Characterizing its order would provide a basis for comparing the precision gains from additional records in either sample, subject to their respective recruitment costs. The central question is the joint behavior of the criterion across sample-size and strength regimes.

\begin{remarkv}{oeq:unequal-sample-frontier}[Unequal-sample length phases]
What is the exact piecewise order of \(R_{N_{\mathrm S},N_{\mathrm T}}(a)\), the unequal-sample minimax expected-length frontier, when the source sample size \(N_{\mathrm S}\) and target sample size \(N_{\mathrm T}\) diverge at unrelated rates? Can a single finite-sample honest interval attain this order in every phase, together with a matching legal instrumental-variable lower family satisfying the no-defiers restriction?
\end{remarkv}

Resolving \cref{obj:oeq:unequal-sample-frontier} would connect an attaining procedure and an information bound for each allocation regime. Requiring the lower family to satisfy the no-defiers restriction keeps that comparison within the same causal model as the proposed intervals.

\subsection{Limitations and future work}

The established benchmark uses a scalar covariate, bounded outcomes, fixed Hölder smoothness with exponent \(\beta=1/8\), and independent records within and across the two samples. Its matching expected-length bounds concern equal source and target sample sizes. The unequal-allocation problem in \cref{obj:def:unequal-sample-handle} defines an extension of the criterion; its general rate characterization and an interval attaining that characterization remain the questions in \cref{obj:oeq:unequal-sample-frontier}.

Further work could examine higher-dimensional covariates, inference under unknown smoothness, and sampling structures with dependent records. These extensions require corresponding coverage and expected-length analyses. Applications with alternative population transport restrictions likewise require an identification analysis suited to those restrictions before the precision benchmark can be extended.
